\documentclass[11pt]{article}
\usepackage[a4paper,margin=25mm]{geometry}
\usepackage[T1]{fontenc}
\usepackage{lmodern}
\usepackage{amsmath,amssymb,amsthm,microtype}
\usepackage[colorlinks=true,linkcolor=blue,citecolor=blue,urlcolor=blue]{hyperref}
\newtheorem{theorem}{Theorem}
\newtheorem{lemma}[theorem]{Lemma}
\newcommand{\N}{\mathbb N}
\newcommand{\lcm}{\operatorname{lcm}}
\hypersetup{pdftitle={The optimal 4/3 density bound in Pilz's problem},pdfauthor={}}
\title{The optimal \(4/3\) density bound in Pilz's problem}
\author{}
\date{}
\begin{document}
\maketitle
\begin{abstract}
For a finite set \(A\) of positive integers, let \(A*[n]\) be the set of
integers with an odd number of representations \(aj\), \(a\in A\),
\(1\le j\le n\). Holdridge and Pach showed that \(|A*[n]|/n\) tends to a
limit \(h(A)\), that \(h(A)\ge8/7\) unless \(A=\{a\}\) or \(A=\{a,2a\}\),
and expected that \(8/7\) could be replaced by \(4/3\). We prove this.
The constant \(4/3\) is best possible, as their computation of
\(h(\{1,3\})\) already shows. The proof is elementary, uses one
divisibility-cut lemma and a calculation for sets of powers of two, and
requires no computer search. It is independent of the resolution of the
finite-interval conjecture in~\cite{Main}. All results are formally
verified in Lean~4.
\end{abstract}

\section*{Introduction}
Here \(\N=\{1,2,\ldots\}\) and \([n]=\{1,\ldots,n\}\), with \([0]=\varnothing\).
For finite \(A,B\subset\N\) we follow~\cite{HP} and write \(A*B\) for the
set of integers with an odd number of ordered representations \(ab\), where
\(a\in A\), \(b\in B\). Thus \(A*[n]\) is the symmetric difference of the
sets \(a[n]\), \(a\in A\).

Pilz's Extended 1--2--3 Conjecture~\cite{Pilz,HKP} asserts that
\(|A*[n]|\ge n\) for every finite nonempty \(A\) and every \(n\ge1\).
Holdridge and Pach~\cite{HP} proved it for all sufficiently large \(n\), and
a proof for every \(n\) is given in~\cite{Main}. Alongside the
finite-interval problem, \cite{HP} studies the asymptotic behaviour for a
fixed set \(A\). Using the exact counting formula of Ke and
Meyer~\cite[(2.1)]{KM}, they show \cite[Lemma 1.6]{HP} that the density
\[
h(A)=\lim_{n\to\infty}\frac{|A*[n]|}{n}
\]
exists, and they prove \cite[Theorem 1.8]{HP} that \(h(A)\ge8/7\) unless
\(A\) has the form \(\{a\}\) or \(\{a,2a\}\). In \cite[Remark 1.9]{HP} they
write that they expect the same theorem to hold with \(8/7\) replaced by
\(4/3\), and that their method could achieve this up to checking a number
of finite cases too large to carry out. They also computed
\(h(\{a,b\})=2-2/b\) for coprime \(a<b\), and showed by a computer search
that the minimum of \(h\) over three-element sets is \(4/3\) and over
four-element sets is \(5/3\)~\cite[Lemma 5.6]{HP}. We prove the bound
\(4/3\) for all sets, without any case check.

\begin{theorem}\label{main}
For every finite nonempty \(A\subset\N\), the limit \(h(A)\) exists and is
rational. If \(A\) is neither \(\{a\}\) nor \(\{a,2a\}\) for any \(a\in\N\),
then \(h(A)\ge4/3\). The excluded sets have density one, and
\(h(\{1,3\})=4/3\), so the constant is best possible.
\end{theorem}

The last sentence is included for completeness; it is already contained
in~\cite[Lemma 5.6]{HP}. The new content is the lower bound \(4/3\). For sets of powers of two the proof
gives the stronger bound \(3/2\).

The argument has two ingredients. The first is a divisibility cut
(Lemma~\ref{cut}), which splits \(A\) according to divisibility by a prime
\(p\) and loses at most \(2/p\) times the density of the part coprime to
\(p\). It is a variant of the splitting estimates of
\cite[Lemma 4.1, Corollary 4.3]{HP}, in which the loss is controlled by a
density rather than by a count of pairs. The second is an explicit
calculation for sets of powers of two, which, after division by the gcd,
are the only sets the cut at an odd prime does not reach.

We first establish the two counting facts needed for the proof.

\begin{lemma}[Existence and scaling]\label{exist}
For every finite nonempty \(A\subset\N\),
\[
h(A)=\sum_{\varnothing\ne J\subseteq A}
       (-2)^{|J|-1}\frac{\min J}{\lcm(J)}.
\]
In particular \(h(A)\) exists and is rational.
Moreover, \(h(cA)=h(A)\) for every \(c\in\N\).
\end{lemma}
The counting formula below is \cite[(2.1)]{KM}, and the limit formula is
\cite[Lemma 1.6]{HP}; we include the short proof for completeness.
\begin{proof}
An integer lying in exactly \(r\) of a finite family of sets contributes
\[
\sum_{j=1}^{r}(-2)^{j-1}\binom rj
=\frac{1-(-1)^r}{2}
\]
to their signed intersection sum. This is one when \(r\) is odd and zero
when \(r\) is even. The intersection of \(a[n]\), \(a\in J\ne\varnothing\),
consists of the positive multiples of \(\lcm(J)\) at most \((\min J)n\).
Consequently,
\[
|A*[n]|=\sum_{\varnothing\ne J\subseteq A}
 (-2)^{|J|-1}\left\lfloor\frac{(\min J)n}{\lcm(J)}\right\rfloor .
\]
Divide by \(n\) and take the limit term by term in this finite sum.
Finally, multiplication by \(c\) preserves all representation counts and
is injective, so \((cA)*[n]=c(A*[n])\).
\end{proof}

\begin{lemma}[The divisibility cut]\label{cut}
Let \(p\) be a prime and let \(A=B\sqcup C\), where no element of \(B\)
is divisible by \(p\) and every element of \(C\) is divisible by \(p\).
Then
\[
h(A)\ \ge\ (1-2/p)h(B)+h(C).
\]
Here the empty set may be assigned density zero.
\end{lemma}
\begin{proof}
The essential identity is
\[
(B*[n])\cap p\N=p\bigl(B*[\lfloor n/p\rfloor]\bigr).
\]
Indeed, if \(py=aj\) with \(a\in B\), then \(p\nmid a\) implies \(j=pk\).
Thus \(py=aj\), \(1\le j\le n\), corresponds bijectively to
\(y=ak\), \(1\le k\le\lfloor n/p\rfloor\). In particular, their
representation counts have the same parity.

Put \(S=B*[n]\) and \(T=C*[n]\). Since \(T\subseteq p\N\) and
\(A*[n]=S\mathbin{\triangle}T\),
\[
|A*[n]|=|S|+|T|-2|S\cap T|
\ge |B*[n]|+|C*[n]|-2|B*[\lfloor n/p\rfloor]|.
\]
Divide by \(n\) and take limits. The final term divided by \(n\)
tends to \(h(B)/p\): set \(m=\lfloor n/p\rfloor\), and use
\((|B*[m]|/m)(m/n)\to h(B)/p\).
\end{proof}

\begin{lemma}[A baseline]\label{baseline}
Every finite nonempty \(A\subset\N\) has \(h(A)\ge1\).
\end{lemma}
\begin{proof}
Induct on \(|A|\). A singleton gives exactly \(n\) products, so has density one.
For \(|A|\ge2\), divide all elements by their greatest common divisor;
Lemma~\ref{exist} leaves the density unchanged. The resulting set has gcd one
and contains an element greater than one. Choose a prime \(p\) dividing that
element, and partition the set into \(B=\{a:p\nmid a\}\) and \(C=\{a:p\mid a\}\).
Both parts are nonempty: the second by our choice, the first because the gcd
is one. By induction their densities are at least one. Since \(p\ge2\),
Lemma~\ref{cut} gives
\[
h(A)\ge(1-2/p)\cdot1+1=2-2/p\ge1.
\]
\end{proof}

\begin{proof}[Proof of Theorem~\ref{main}]
Existence and rationality of \(h(A)\) are Lemma~\ref{exist}.
Divide \(A\) by its gcd. This preserves its density and whether it belongs
to either excluded family. If an odd prime \(p\) divides one of its elements,
the partition in the preceding proof and Lemmas~\ref{cut}--\ref{baseline} give
\[
h(A)\ge2-2/p\ge4/3.
\]

Otherwise every element is a power of two. In such a set, deleting the least
element cannot increase the density: divide by that element, leaving
\(B=\{1\}\) and a set \(C\) of even numbers, and apply Lemma~\ref{cut} with
\(p=2\) to obtain \(h(B\cup C)\ge h(C)\). Scaling back proves the assertion.

For two powers of two, scale to \(A=\{1,u\}\). Outside the excluded case
\(u=2\), we have \(u\ge4\), and Lemma~\ref{exist} gives
\[
h(A)=2-\frac2u\ge\frac32.
\]
For three powers of two, scale to \(A=\{1,u,uv\}\), where \(u,v\ge2\).
The same formula gives
\[
h(A)=3-\frac2u-\frac2v+\frac2{uv}
=1+2\left(1-\frac1u\right)\left(1-\frac1v\right)\ge\frac32.
\]
For a larger set, delete the least element repeatedly until three remain.
Its original density is at least that of the final three-element set.
This proves the required bound, and in fact the stronger \(3/2\) bound
in the powers-of-two case.

Finally, a singleton has density one, and direct counting gives
\[
|\{a,qa\}*[n]|=2n-2\lfloor n/q\rfloor\qquad(q\ge2).
\]
Taking \(q=2\) proves the remaining exception; taking \(a=1,q=3\)
proves sharpness.
\end{proof}

\paragraph{Scope.}
The set \(A\) is fixed while \(n\) tends to infinity.
There is no assertion that \(|A*[n]|\ge4n/3\) for every \(n\):
for example, \(\{1,2,3\}*[3]=\{1,4,9\}\).
That density one occurs exactly in the two stated families already
follows from~\cite[Theorem 1.8]{HP}. The proof does not classify all sets
of density \(4/3\); among sets of greatest common divisor one with elements
at most \(12\), the equality cases are \(\{1,3\}\), \(\{2,3\}\) and
\(\{1,2,3\}\).

\section*{Formalisation}
The accompanying Lean~4 sources give a formal proof of Theorem~\ref{main}.
The complete statement is \texttt{PilzDensity.companion\_theorem}; optimality
of \(4/3\) is \texttt{PilzDensity.four\_thirds\_optimal}, and the limit for
\(\{1,3\}\) is \texttt{PilzDensity.sharpness\_limit}. The density is defined
directly as the limit of the ordered odd-representation count, and the
development does not depend on the finite-interval theorem of~\cite{Main}.
The repository's \texttt{paper/README.md} maps each lemma to its Lean
theorem and gives build and verification commands.

\section*{Acknowledgements}
The question answered here was raised by Philippa Holdridge and
P\'eter P\'al Pach~\cite[Remark 1.9]{HP}, whose paper also supplies the
density framework and the splitting estimates on which the divisibility
cut is modelled. The counting formula used throughout is recorded by Ke
and Meyer~\cite[(2.1)]{KM}.

\paragraph{How this work was produced.}
The arguments in this note, its written proofs, and its formalisation were
developed with substantial assistance from AI systems. The written proofs
have not been checked line by line by a human mathematician. Correctness
rests on the Lean~4 formalisation: every stated theorem is machine-checked
and depends only on the standard axioms of Lean's foundations. Where the
text and the formal proof differ, the formal proof is authoritative.

\section*{License}
This paper is licensed under the Creative Commons Attribution~4.0
International License, \url{https://creativecommons.org/licenses/by/4.0/}.

\begin{thebibliography}{9}
\bibitem{HP}
P. Holdridge and P. P. Pach,
\emph{On the Extended 1-2-3 Conjecture of Pilz},
preprint (2026), \href{https://arxiv.org/abs/2607.00934}{arXiv:2607.00934}.
\bibitem{HKP}
P.-Y. Huang, W.-F. Ke and G. F. Pilz,
\emph{The cardinality of some symmetric differences},
Proc. Amer. Math. Soc. \textbf{138} (2010), 787--797.
\bibitem{KM}
W.-F. Ke and J. H. Meyer,
\emph{On a conjecture regarding the symmetric difference of certain sets},
Bull. Aust. Math. Soc. \textbf{111} (2025), 397--404.
\bibitem{Main}
\emph{Pilz's Extended 1--2--3 Conjecture over every field: sparse
multiples of weighted integer intervals}, preprint (2026), with Lean~4
formalisation; an earlier proof of the binary case is in the same archive,
\emph{Pilz's Extended 1--2--3 Conjecture: proofs and Lean~4
formalisations}, Zenodo, 2026,
\href{https://doi.org/10.5281/zenodo.23568997}{doi:10.5281/zenodo.23568997}.
\bibitem{Pilz}
G. F. Pilz,
\emph{On polynomial near-ring codes},
Contrib. Gen. Algebra \textbf{8} (1992), 233--238.
\end{thebibliography}
\end{document}
